\chapter{Computing with Matrices}\label{ch:matrices}

\begin{center}\itshape
I felt sure that none of the variables could have more than six correct figures and yet the agreement was as good\\
as it would have been if I had been given the exact answer and had then rounded it to ten figures.\\
\upshape(James H. Wilkinson, Turing Award lecture, 1970)
\end{center}

In 1946 a young mathematician named James Hardy Wilkinson arrived at the National Physical Laboratory near London, to work half of his time with the
computing section of Leslie Fox and Charles Goodwin and the other half with Alan Turing, who was designing an electronic computer. He found the
laboratory in a gloomy mood. Bounds had been published which suggested that the error in the solution of a system of $n$ linear equations by
elimination could grow like $4^n$, and as Wilkinson remembered it, ``it was the more distinguished mathematicians who were most pessimistic''. If the
bounds were right, the new machines would be useless for systems of quite modest size.

Wilkinson had a secret. During the war, at the Armament Research Department, he had been handed twelve equations in twelve unknowns and had
solved them by hand on a desk calculator, by Gaussian elimination. The last reduced equation had a coefficient of order $10^{-5}$, so he expected at most six
correct figures, and he computed everything to ten figures only because a fixed routine causes fewer blunders. Then he substituted his
answer back into the twelfth equation, and the two sides agreed to ten figures. ``That, I said to myself, was a coincidence.'' Eleven more
coincidences followed. Not wanting to look foolish, he kept quiet.

Then a system of eighteen equations arrived at the laboratory, and Fox, Goodwin, Turing and Wilkinson decided to solve it by hand, as a team.
Turing ``was convinced that it would be a failure''. History repeated itself: the residuals were of order $10^{-10}$, as small as for the exact
solution rounded to ten decimals, while a later check showed the solution itself had only about six correct figures. ``I suppose this must be
regarded as a defeat for Turing'', Wilkinson wrote, but ``the second round undoubtedly went to Turing'': a year later Turing published
\emph{Rounding-off errors in matrix processes}, which, together with a long paper of John von Neumann and Herman Goldstine, began to dispel the gloom.
It was Turing who named the \emph{condition number}.

The resolution of the puzzle, worked out by Wilkinson over the next fifteen years, is the subject of this chapter. Its key idea is to stop asking
how wrong the computed answer is, and to ask instead: for which nearby problem is it exactly right? Wilkinson's computed solution was the exact
solution of a system differing from his in the tenth figure; whether that makes the answer accurate depends on the system, not on the algorithm, and is
measured by the condition number. We prove this for triangular systems, Gaussian elimination and the $QR$ factorisation by Householder reflections,
then turn to the harder problem of eigenvalues, where the $QR$ algorithm, invented by John Francis in 1959--1961, is one of the great algorithms of
the century, and finally to the singular value decomposition. The theory of the SVD, of the pseudo-inverse and of least squares was done in Tome~S1,
Chapter~6; this chapter is its numerical side, as promised there.

\begin{think}
The Hilbert matrix $H_n=\bigl(1/(i+j-1)\bigr)_{1\le i,j\le n}$ is symmetric positive definite. Take $x=(1,\dots,1)$, compute $b=H_{10}x$, and solve
$H_{10}y=b$ with a good library routine, working with $16$ significant digits. How many correct digits do you expect in $y$? And how small do you expect
the residual $H_{10}y-b$ to be?
\end{think}

\section{Floating-point arithmetic}\label{sec:float}

A computer represents real numbers by a finite set $\mathbb F\subset\R$ of \emph{floating-point numbers} $\pm m\cdot2^{e}$, where the
\emph{significand} $m$ has $t$ binary digits, $2^{t-1}\le m<2^t$, and the exponent $e$ lies in a range wide enough that we may ignore overflow and
underflow. In the IEEE~754 standard of 1985 used by essentially every processor, double precision has $t=53$. Every real number $x$ in the range
has a nearest element $\mathrm{fl}(x)$ of $\mathbb F$, and
\[
\mathrm{fl}(x)=x(1+\delta),\qquad|\delta|\le u=2^{-t},
\]
where $u$ is the \emph{unit roundoff}: $u=2^{-53}\approx1.11\times10^{-16}$ in double precision. Indeed, if $2^{e+t-1}\le|x|<2^{e+t}$, consecutive
floating-point numbers near $x$ are $2^e$ apart, so $|\mathrm{fl}(x)-x|\le2^{e-1}\le2^{-t}|x|$. The standard requires the four operations and the
square root to be \emph{correctly rounded}: the computed result is the rounded exact result. This gives the only property we shall use.

\begin{definition}[standard model]\label{def:model}
For $a,b\in\mathbb F$ and $\mathrm{op}\in\{+,-,\times,/\}$,
\[
\mathrm{fl}(a\mathbin{\mathrm{op}}b)=(a\mathbin{\mathrm{op}}b)(1+\delta),\qquad\mathrm{fl}(\sqrt a)=\sqrt a\,(1+\delta),\qquad|\delta|\le u.
\]
\end{definition}

So each single operation is almost perfect, with relative error at most $10^{-16}$. The danger is in what follows: errors accumulate, and above all
they are amplified when nearly equal numbers are subtracted. In $\mathrm{fl}(a-b)$ with $a=1+10^{-12}$ and $b=1$, both computed earlier with relative
errors of order $u$, the subtraction is exact, but the absolute errors of $a$ and $b$, of order $10^{-16}$, now sit on a result of size $10^{-12}$: the
relative error has become $10^{-4}$. This is \emph{cancellation}, and it is never the fault of the subtraction, only of what came before.

Errors that look negligible can also add up over a long run. On 25 February 1991 a Patriot missile battery at Dhahran, Saudi Arabia, failed to
intercept an incoming Scud missile, which hit an army barracks and killed 28 American soldiers. The system counted time in tenths of a second and
multiplied by $1/10$ in a 24-bit fixed-point register. The binary expansion of $1/10=0.000110011001100\dots_2$ does not terminate, and chopping it after
23 binary places makes an error of about $9.5\times10^{-8}$ per tenth of a second. After the 100 hours the battery had been running, $3.6$ million
tenths, the clock was off by $0.3433$ seconds, during which a Scud travels more than half a kilometre; the government report lists a shift of 687
metres in the range gate where the radar looked for the target. A corrected version of the software reached Dhahran the next day.

We shall carry products of many factors $1+\delta_i$, and bound them by the following device.

\begin{lemma}[accumulating roundoff]\label{lem:gamma}
For $k\ge1$ with $ku<1$ put $\gamma_k=ku/(1-ku)$.
\begin{enumerate}
\item If $|\delta_i|\le u$ and $\rho_i\in\{1,-1\}$ for $i\le k$, then $\prod_{i=1}^k(1+\delta_i)^{\rho_i}=1+\theta$ with $|\theta|\le\gamma_k$.
\item $\gamma_j\le\gamma_k$ for $j\le k$, $c\gamma_k\le\gamma_{ck}$ for integers $c\ge1$, and $\gamma_j+\gamma_k+\gamma_j\gamma_k\le\gamma_{j+k}$.
\end{enumerate}
\end{lemma}
\begin{proof}
(1) Each factor lies in $[1-u,(1-u)^{-1}]$, since $1+u\le(1-u)^{-1}$ and $(1+u)^{-1}\ge1-u$. Hence the product lies between $(1-u)^k\ge1-ku\ge1-\gamma_k$
(Bernoulli's inequality) and $(1-u)^{-k}\le(1-ku)^{-1}=1+\gamma_k$. (2) The first two follow from $ck u/(1-ku)\le cku/(1-cku)$. For the third,
$1+\gamma_j+\gamma_k+\gamma_j\gamma_k=(1+\gamma_j)(1+\gamma_k)=\frac1{(1-ju)(1-ku)}\le\frac1{1-(j+k)u}=1+\gamma_{j+k}$, as $(1-ju)(1-ku)\ge1-(j+k)u$.
\end{proof}

In double precision $\gamma_k$ is $ku$ up to a factor $1.0001$ as long as $k\le10^{12}$; think of $\gamma_k$ as ``$k$ roundings''. We write $|A|$ for
the matrix of absolute values of the entries of $A$, and inequalities between matrices or vectors entrywise.

\begin{proposition}[inner products]\label{prop:dot}
For $x,y\in\mathbb F^n$, the inner product computed by recursive summation $s_1=x_1y_1$, $s_i=s_{i-1}+x_iy_i$ satisfies
\[
\mathrm{fl}(x^{\mathsf T}y)=\sum_{i=1}^nx_iy_i(1+\theta_i),\quad|\theta_i|\le\gamma_n,\qquad\text{so}\quad\bigl|\mathrm{fl}(x^{\mathsf T}y)-x^{\mathsf T}y\bigr|\le\gamma_n\,|x|^{\mathsf T}|y|.
\]
Consequently $\mathrm{fl}(Ax)=(A+\Delta A)x$ for a matrix $A$ with $n$ columns, with $|\Delta A|\le\gamma_n|A|$.
\end{proposition}
\begin{proof}
The term $x_iy_i$ is rounded once when formed and once in each of the additions $s_i,\dots,s_n$ (the first term in $s_2,\dots,s_n$): at most $n$
factors $1+\delta$, and Lemma~\ref{lem:gamma} applies. Each row of $\mathrm{fl}(Ax)$ is such an inner product, $\sum_ja_{ij}(1+\theta_{ij})x_j$, and we put
$(\Delta A)_{ij}=a_{ij}\theta_{ij}$.
\end{proof}

The second statement is the model of everything that follows: the computed product is the \emph{exact} product of a slightly perturbed matrix. When
$x^{\mathsf T}y$ is much smaller than $|x|^{\mathsf T}|y|$, that is, when cancellation occurs, the computed inner product can still have a large
relative error; but it is exactly right for slightly different data.

\section{Conditioning and backward stability}\label{sec:conditioning}

Let us separate two questions. How sensitive is the \emph{problem} to its data? And how good is the \emph{algorithm} at solving it? We use the
Euclidean norm on vectors and the induced operator norm $\|A\|=\max_{\|x\|=1}\|Ax\|=\sigma_1(A)$ on matrices (Tome~S1, Chapter~6), and occasionally the
Frobenius norm $\|A\|_F=(\sum a_{ij}^2)^{1/2}$.

\begin{definition}\label{def:cond}
The \emph{condition number} of an invertible matrix $A$ is $\kappa(A)=\|A\|\,\|A^{-1}\|=\sigma_1(A)/\sigma_n(A)\ge1$.
\end{definition}

\begin{theorem}[perturbation of linear systems]\label{thm:perturb}
Let $Ax=b$ with $A$ invertible and $b\ne0$, and $(A+\Delta A)(x+\Delta x)=b+\Delta b$ with $\|\Delta A\|\le\varepsilon\|A\|$, $\|\Delta b\|\le\varepsilon\|b\|$
and $\varepsilon\kappa(A)<1$. Then
\[
\frac{\|\Delta x\|}{\|x\|}\le\frac{2\varepsilon\kappa(A)}{1-\varepsilon\kappa(A)}.
\]
\end{theorem}
\begin{proof}
Subtracting $Ax=b$ gives $(A+\Delta A)\Delta x=\Delta b-\Delta A\,x$, so $\Delta x=A^{-1}(\Delta b-\Delta A\,x-\Delta A\,\Delta x)$ and
\[
\|\Delta x\|\le\|A^{-1}\|\bigl(\varepsilon\|b\|+\varepsilon\|A\|\|x\|+\varepsilon\|A\|\|\Delta x\|\bigr).
\]
Since $\|b\|\le\|A\|\|x\|$, this is $\|\Delta x\|\le2\varepsilon\kappa(A)\|x\|+\varepsilon\kappa(A)\|\Delta x\|$, which rearranges to the claim.
\end{proof}

The bound is attained up to the factor $2$: take $b$ along the top left singular vector $u_1$ and $\Delta b=\varepsilon\|b\|u_n$, so that
$\|x\|=\|b\|/\sigma_1$ and $\|\Delta x\|=\varepsilon\|b\|/\sigma_n$. So $\kappa(A)$ is the factor by which relative errors in the data may be amplified, and
$\log_{10}\kappa(A)$ the number of decimal digits one should expect to lose, \emph{whatever the algorithm}. The condition number has a second meaning.

\begin{theorem}[distance to singularity; Gastinel, Kahan]\label{thm:distsing}
For invertible $A$, $\min\{\|\Delta A\|/\|A\|:A+\Delta A\text{ singular}\}=1/\kappa(A)$.
\end{theorem}
\begin{proof}
If $A+\Delta A$ is singular, take $x\ne0$ with $Ax=-\Delta A\,x$; then $\|x\|=\|A^{-1}\Delta A\,x\|\le\|A^{-1}\|\|\Delta A\|\|x\|$, so
$\|\Delta A\|\ge1/\|A^{-1}\|=\sigma_n$. Conversely, with the SVD $A=\sum_i\sigma_iu_iv_i^{\mathsf T}$, the matrix $\Delta A=-\sigma_nu_nv_n^{\mathsf T}$ has norm
$\sigma_n$ and $(A+\Delta A)v_n=0$. Divide by $\|A\|=\sigma_1$.
\end{proof}

An ill-conditioned matrix is one that is relatively close to a singular one. Hilbert matrices are the classical example: $\kappa(H_{10})\approx1.6\times10^{13}$
and $\kappa(H_{12})\approx1.6\times10^{16}$, so $H_{12}$ is within rounding distance of a singular matrix (Exercise~\ref{ex:hilbertcond} explains the
growth, about $e^{3.5n}$).

\begin{definition}[backward stability]\label{def:backward}
An algorithm computing $y=f(x)$ is \emph{backward stable} if, for every input $x$, the computed $\hat y$ equals $f(x+\Delta x)$ exactly for some
$\Delta x$ with $\|\Delta x\|\le Cu\|x\|$, where $C$ is a modest function of the dimensions.
\end{definition}

Wilkinson's riddle is now resolved in one line. A backward stable algorithm for $Ax=b$ returns the exact solution of a system whose data differ from
$(A,b)$ in the last digits, so by Theorem~\ref{thm:perturb} its relative error is at most about $2Cu\kappa(A)$: forward error $\lesssim$ condition number
$\times$ backward error. And its residual is tiny: $\|b-A\hat x\|=\|\Delta A\,\hat x-\Delta b\|\le Cu(\|A\|\|\hat x\|+\|b\|)$. With $\kappa$ of order $10^4$ or $10^5$ and
ten figures, Wilkinson's residuals had to be of order $10^{-10}$ and his solutions could have five or six correct figures, as he observed. For $H_{10}$
in double precision, expect a residual of order $10^{-16}$ and at least $16-13=3$ correct digits; the notebook finds about five, since the bound is a worst case.

\section{Triangular systems and Gaussian elimination}\label{sec:ge}

All the elimination algorithms repeatedly solve one linear equation for one unknown, and their analysis rests on the following lemma.

\begin{lemma}[one equation]\label{lem:oneeq}
Let $y=(c-\sum_{k=1}^{s-1}a_kb_k)/b_s$ be computed as $c$ minus the products $a_kb_k$ subtracted one at a time, followed by one division (or no
division if $b_s=1$). Then the computed $\hat y$ satisfies
\[
b_s\hat y(1+\theta_s)=c-\sum_{k=1}^{s-1}a_kb_k(1+\theta_k),\qquad|\theta_k|\le\gamma_s\quad(1\le k\le s).
\]
\end{lemma}
\begin{proof}
Put $t_0=c$ and $t_k=\bigl(t_{k-1}-a_kb_k(1+\delta_k)\bigr)(1+\delta_k')$, so $\hat y=t_{s-1}(1+\delta'')/b_s$. Unrolling,
$t_{s-1}=c\prod_{k}(1+\delta_k')-\sum_ka_kb_k(1+\delta_k)\prod_{l\ge k}(1+\delta'_l)$. Divide by $\prod_k(1+\delta'_k)$:
\[
b_s\hat y\,(1+\delta'')^{-1}\prod_{k=1}^{s-1}(1+\delta_k')^{-1}=c-\sum_{k=1}^{s-1}a_kb_k(1+\delta_k)\prod_{l<k}(1+\delta_l')^{-1}.
\]
The left factor has $s$ roundings, the $k$th term on the right at most $k\le s-1$; apply Lemma~\ref{lem:gamma}.
\end{proof}

\begin{theorem}[back substitution is backward stable]\label{thm:backsub}
Let $T$ be an invertible upper triangular $n\times n$ matrix, and let $\hat x$ be computed by $x_n=b_n/t_{nn}$ and
$x_i=(b_i-\sum_{j>i}t_{ij}x_j)/t_{ii}$ for $i=n-1,\dots,1$. Then $(T+\Delta T)\hat x=b$ with $|\Delta T|\le\gamma_n|T|$.
\end{theorem}
\begin{proof}
Row $i$ is Lemma~\ref{lem:oneeq} with $s=n-i+1$: $t_{ii}\hat x_i(1+\theta_{ii})+\sum_{j>i}t_{ij}\hat x_j(1+\theta_{ij})=b_i$, with all
$|\theta_{ij}|\le\gamma_{n-i+1}\le\gamma_n$. Put $(\Delta T)_{ij}=t_{ij}\theta_{ij}$.
\end{proof}

Note that $\Delta T$ depends on $b$; that is allowed, since backward stability is a statement about one input at a time.

Gaussian elimination reduces $Ax=b$ to triangular form. At step $k$ it subtracts $l_{ik}=a^{(k)}_{ik}/a^{(k)}_{kk}$ times row $k$ from row $i>k$; the
multipliers form a unit lower triangular $L$, the final rows an upper triangular $U$, and $A=LU$. With \emph{partial pivoting} one first swaps row $k$
with the row $i\ge k$ having the largest $|a^{(k)}_{ik}|$, so that $|l_{ik}|\le1$; this amounts to $PA=LU$ for a permutation matrix $P$, which we absorb
into $A$. Without pivoting the method can fail stupidly: for $A=\begin{psmallmatrix}\varepsilon&1\\1&1\end{psmallmatrix}$ with $\varepsilon=10^{-20}$, the
multiplier is $10^{20}$, the computed $u_{22}=\mathrm{fl}(1-10^{20})=-10^{20}$ forgets the $1$, and the computed solution of $Ax=(1,2)$ is $(0,1)$ instead
of approximately $(1,1)$ (Exercise~\ref{ex:nopivot}).

\begin{theorem}[Wilkinson, 1961]\label{thm:geerror}
Gaussian elimination (with the rows already in pivot order) computes factors $\hat L,\hat U$ with $\hat L\hat U=A+\Delta A$, $|\Delta A|\le\gamma_n|\hat L||\hat U|$.
Solving $Ax=b$ by $\hat Ly=b$ and $\hat Ux=\hat y$ gives $(A+E)\hat x=b$ with $|E|\le\gamma_{3n}|\hat L||\hat U|$.
\end{theorem}
\begin{proof}
The entry $a_{ij}$, $i\le j$, is transformed in the steps $k=1,\dots,i-1$ by $a_{ij}^{(k+1)}=a_{ij}^{(k)}-\hat l_{ik}\hat u_{kj}$, and the last value is
$\hat u_{ij}$: this is the computation of Lemma~\ref{lem:oneeq} with $c=a_{ij}$, $s=i$ and $b_s=1$ (no division), giving
$\hat u_{ij}(1+\theta_{ii})=a_{ij}-\sum_{k<i}\hat l_{ik}\hat u_{kj}(1+\theta_k)$. For $i>j$, the entry is transformed in the steps $k<j$ and then divided by
$\hat u_{jj}$ to give $\hat l_{ij}$: Lemma~\ref{lem:oneeq} with $s=j$ and $b_s=\hat u_{jj}$. In both cases $a_{ij}=\sum_k\hat l_{ik}\hat u_{kj}(1+\theta'_{k})$
with $|\theta'_k|\le\gamma_n$, since $\hat l_{ii}=1$; this is the first claim. By Theorem~\ref{thm:backsub} and its lower triangular twin,
$(\hat L+\Delta L)\hat y=b$ and $(\hat U+\Delta U)\hat x=\hat y$ with $|\Delta L|\le\gamma_n|\hat L|$, $|\Delta U|\le\gamma_n|\hat U|$. So
$b=(\hat L+\Delta L)(\hat U+\Delta U)\hat x=(A+E)\hat x$ with $E=\Delta A+\Delta L\hat U+\hat L\Delta U+\Delta L\Delta U$, and
$|E|\le(3\gamma_n+\gamma_n^2)|\hat L||\hat U|\le\gamma_{3n}|\hat L||\hat U|$ by Lemma~\ref{lem:gamma}.
\end{proof}

So elimination is backward stable \emph{provided $|\hat L||\hat U|$ is not much larger than $|A|$}. With partial pivoting $|\hat l_{ij}|\le1$, and everything
depends on the \emph{growth factor} $\rho_n=\max_{i,j,k}|a^{(k)}_{ij}|/\max_{i,j}|a_{ij}|$, since then
$\||\hat L||\hat U|\|_\infty\le n^2\rho_n\max|a_{ij}|$.

\begin{proposition}[growth]\label{prop:growth}
With partial pivoting $\rho_n\le2^{n-1}$, and this bound is attained.
\end{proposition}
\begin{proof}
Since $|l_{ik}|\le1$, each step gives $|a_{ij}^{(k+1)}|\le|a_{ij}^{(k)}|+|a_{kj}^{(k)}|\le2\max|a^{(k)}|$, and there are $n-1$ steps. For the example, let $W_n$ have
$1$ on the diagonal and in the last column, $-1$ below the diagonal, and $0$ elsewhere. No row swaps occur (the pivots and all entries below them have modulus
$1$, and ties keep the current row), and step $k$ adds row $k$ to every later row: the last column doubles at each step, ending with $u_{nn}=2^{n-1}$
(Exercise~\ref{ex:wilkinsongrowth}).
\end{proof}

So Gaussian elimination with partial pivoting is not backward stable in the worst case. Yet it is the method every library uses, because matrices
with large growth are extraordinarily rare. In 1990 Lloyd Trefethen and Robert Schreiber reported that for random matrices the growth factor
stays small, growing only like a small power of $n$ (about $\sqrt n$ in the notebook's experiments), and in 2006 Arvind Sankar, Daniel Spielman and Shang-Hua Teng proved, in the spirit of smoothed analysis, that small random
perturbations of any matrix make large growth unlikely (for elimination without pivoting). Why the bad matrices are so rare is still not fully
understood. The next algorithm has no such mystery.

\section{Householder reflections and the QR factorisation}\label{sec:householder}

Orthogonal matrices preserve lengths, so multiplying by them cannot amplify errors. In 1958 Alston Householder proposed to triangularise a matrix by
reflections. For $v\ne0$, the \emph{Householder reflection}
\[
P=I-\frac{2vv^{\mathsf T}}{v^{\mathsf T}v}
\]
is the reflection in the hyperplane $v^\perp$: $Pv=-v$ and $Pw=w$ for $w\perp v$; so $P$ is symmetric, orthogonal and $P^2=I$. Given $x\ne0$, the
reflection exchanging $x$ and $y$ with $\|y\|=\|x\|$ is the one with $v=x-y$ (Figure~\ref{fig:householder}). To send $x$ to a multiple of $e_1$, the two
choices are $y=\pm\|x\|e_1$, and we choose
\[
v=x+\operatorname{sign}(x_1)\|x\|e_1,\qquad Px=-\operatorname{sign}(x_1)\|x\|e_1\qquad(\operatorname{sign}0=1).
\]
The sign is chosen so that $v_1=x_1+\operatorname{sign}(x_1)\|x\|$ adds two numbers of the same sign: no cancellation, even when $x$ is nearly parallel
to $e_1$. That every orthogonal matrix of $\R^n$ is a product of at most $n$ such reflections is the Cartan--Dieudonn\'e theorem (Tome~S1, Exercise~4.13).

\begin{figure}[ht]
\centering
\fig{ch01-householder}
\caption{The Householder reflection with $v=x+\|x\|e_1$ maps $x$ to $-\|x\|e_1$, mirroring it in the hyperplane $v^\perp$. The other choice, $v'=x-\|x\|e_1$,
would be short when $x$ is close to the $e_1$ axis and would be computed with cancellation.}
\label{fig:householder}
\end{figure}

\emph{Householder $QR$.} Let $A$ be $m\times n$ with $m\ge n$. Step $k=1,\dots,n$ takes the part $x=(a_{kk},\dots,a_{mk})$ of column $k$ on and below the
diagonal, forms $v_k$ from it, and applies $P_k=I_{k-1}\oplus(I-2v_kv_k^{\mathsf T}/v_k^{\mathsf T}v_k)$ to the columns $k+1,\dots,n$, while column $k$ is
replaced by $(-\operatorname{sign}(x_1)\|x\|,0,\dots,0)$. After $n$ steps $P_n\cdots P_1A=R$ is upper triangular, so $A=QR$ with $Q=P_1\cdots P_n$
orthogonal. One never forms $P_k$: applying it to a vector $b$ costs one inner product and one update, $b-\beta v(v^{\mathsf T}b)$ with
$\beta=2/v^{\mathsf T}v$, about $4(m-k)$ operations; the whole factorisation costs about $2mn^2-\frac23n^3$ operations.

\begin{lemma}[one reflection]\label{lem:onereflection}
Let $x\in\mathbb F^m$, $x\ne0$, and let $\hat v$ be the computed $v$, $\tilde P=I-2\hat v\hat v^{\mathsf T}/\hat v^{\mathsf T}\hat v$ the \emph{exact} reflection
defined by it, and $\eta=\gamma_{10m+20}$. Then:
\begin{enumerate}
\item for every $b\in\mathbb F^m$, the computed $\hat y=\mathrm{fl}(b-\hat\beta\hat v(\hat v^{\mathsf T}b))$ equals $\tilde P(b+\Delta b)$ with $\|\Delta b\|\le\eta\|b\|$;
\item the computed image $\hat r e_1$ of $x$, with $\hat r=-\operatorname{sign}(x_1)\hat\sigma$ and $\hat\sigma=\mathrm{fl}(\|x\|)$, equals
$\tilde P(x+\Delta x)$ with $\|\Delta x\|\le\eta\|x\|$.
\end{enumerate}
\end{lemma}
\begin{proof}
\emph{The vector.} By Proposition~\ref{prop:dot} the computed $\sum x_i^2$ is $\|x\|^2(1+\theta)$ with $|\theta|\le\gamma_m$ (all terms are
nonnegative), and $\hat\sigma=\|x\|(1+\theta)^{1/2}(1+\delta)$, so $|\hat\sigma-\|x\||\le\gamma_{m+1}\|x\|$ (as $|(1+\theta)^{1/2}-1|\le|\theta|$). Then
$\hat v_1=(x_1+\operatorname{sign}(x_1)\hat\sigma)(1+\delta)$, and since $x_1$ and $\operatorname{sign}(x_1)\hat\sigma$ have the same sign,
$|x_1+\operatorname{sign}(x_1)\hat\sigma-v_1|=|\hat\sigma-\|x\||\le\gamma_{m+1}\|x\|\le\gamma_{m+1}|v_1|$. With Lemma~\ref{lem:gamma}, $|\hat v_1-v_1|\le\gamma_{m+2}|v_1|$,
while $\hat v_i=x_i=v_i$ for $i\ge2$. Hence $\|\hat v-v\|\le\gamma_{m+2}\|v\|$.

(1) Let $\tilde\beta=2/\hat v^{\mathsf T}\hat v$. The computed $\hat\beta$ is $\tilde\beta(1+\theta')$ with $|\theta'|\le\gamma_{m+1}$. The computed
$w=\mathrm{fl}(\hat v^{\mathsf T}b)$ is $\hat v^{\mathsf T}b+e$ with $|e|\le\gamma_m\|\hat v\|\|b\|$ (Proposition~\ref{prop:dot} and Cauchy--Schwarz). Then
$c=\mathrm{fl}(\hat\beta w)$ and $z_i=\mathrm{fl}(c\hat v_i)$ give $z_i=\tilde\beta(\hat v^{\mathsf T}b+e)(1+\theta'')\hat v_i$ with $|\theta''|\le\gamma_{m+3}$, so
\[
\|z-\tilde\beta(\hat v^{\mathsf T}b)\hat v\|\le\tilde\beta\|\hat v\|\bigl(\gamma_{m+3}\|\hat v\|\|b\|+(1+\gamma_{m+3})\gamma_m\|\hat v\|\|b\|\bigr)\le2\gamma_{2m+3}\|b\|,
\]
using $\tilde\beta\|\hat v\|^2=2$ and Lemma~\ref{lem:gamma}. Finally $\hat y_i=(b_i-z_i)(1+\delta_i)$, so $\hat y-\tilde Pb=(\tilde\beta(\hat v^{\mathsf T}b)\hat v-z)+\delta\circ(b-z)$,
where $\|z\|\le\|\tilde\beta(\hat v^{\mathsf T}b)\hat v\|+2\gamma_{2m+3}\|b\|\le(2+2\gamma_{2m+3})\|b\|$. Altogether
$\|\hat y-\tilde Pb\|\le(2\gamma_{2m+3}+3u+2u\gamma_{2m+3})\|b\|\le\gamma_{4m+9}\|b\|$ by Lemma~\ref{lem:gamma}. Put $\Delta b=\tilde P(\hat y-\tilde Pb)$; then
$\tilde P(b+\Delta b)=\tilde Pb+(\hat y-\tilde Pb)=\hat y$, as $\tilde P^2=I$, and $\|\Delta b\|=\|\hat y-\tilde Pb\|$.

(2) Let $P$ be the exact reflection defined by $v$, so $Px=re_1$ with $r=-\operatorname{sign}(x_1)\|x\|$. For unit vectors $p,q$,
$\|pp^{\mathsf T}-qq^{\mathsf T}\|=\|p(p-q)^{\mathsf T}+(p-q)q^{\mathsf T}\|\le2\|p-q\|$, and for nonzero $a,b$,
$\bigl\|\frac a{\|a\|}-\frac b{\|b\|}\bigr\|\le\frac{\|a-b\|}{\|b\|}+\frac{|\|b\|-\|a\||}{\|b\|}\le\frac{2\|a-b\|}{\|b\|}$. Hence
$\|\tilde P-P\|\le4\cdot2\gamma_{m+2}=8\gamma_{m+2}$. Put $\Delta x=\tilde P(\hat re_1)-x$. Then $\tilde P(x+\Delta x)=\hat re_1$ and
\[
\|\Delta x\|=\|\hat re_1-\tilde Px\|\le|\hat r-r|+\|Px-\tilde Px\|\le(\gamma_{m+1}+8\gamma_{m+2})\|x\|\le\gamma_{9m+17}\|x\|.
\]
Both bounds are at most $\eta=\gamma_{10m+20}$.
\end{proof}

We made no effort to optimise the constant $10m+20$; a finer analysis gives about $m$, and what matters is that it is a modest multiple of $u$.

\begin{theorem}[Householder $QR$ is backward stable]\label{thm:qrstable}
Let $\hat R$ be the triangular factor computed by Householder $QR$ of an $m\times n$ matrix $A$ in floating point, and $\eta=\gamma_{10m+20}$ with
$n\eta<1$. There is an exactly orthogonal $Q$ with
\[
A+\Delta A=Q\hat R,\qquad\|\Delta a_j\|\le\frac{n\eta}{1-n\eta}\|a_j\|\quad\text{for each column }j.
\]
Moreover, the computed $\widehat{Q^{\mathsf T}b}$ for a vector $b$, obtained by applying the computed reflections in turn, equals $Q^{\mathsf T}(b+\Delta b)$ with
$\|\Delta b\|\le\frac{n\eta}{1-n\eta}\|b\|$, for the same $Q$.
\end{theorem}
\begin{proof}
Let $\tilde P_k=I_{k-1}\oplus\tilde P$ be the exact reflection of step $k$ defined by the computed $\hat v_k$, and $Q=\tilde P_1\cdots\tilde P_n$. Fix a column
$j$ and let $a^{(k)}$ be its computed value after step $k$, with $a^{(0)}=a_j$. At each step $k\le j$, the active part of $a^{(k-1)}$ is transformed as in
Lemma~\ref{lem:onereflection} (part (1) if $k<j$, part (2) if $k=j$), while its first $k-1$ entries are untouched; so
$a^{(k)}=\tilde P_k(a^{(k-1)}+\Delta_k)$ with $\|\Delta_k\|\le\eta\|a^{(k-1)}\|$, and $\|a^{(k)}\|\le(1+\eta)\|a^{(k-1)}\|$. Steps $k>j$ leave the column alone,
since its entries in rows $k,\dots,m$ are zero. Multiplying by $\tilde P_1\cdots\tilde P_k$ and using $\tilde P_i^2=I$,
\[
\tilde P_1\cdots\tilde P_ka^{(k)}=\tilde P_1\cdots\tilde P_{k-1}a^{(k-1)}+\tilde P_1\cdots\tilde P_{k-1}\Delta_k,
\]
and telescoping, $Q\hat r_j=a_j+\Delta a_j$ with $\Delta a_j=\sum_{k\le j}\tilde P_1\cdots\tilde P_{k-1}\Delta_k$. Hence
\[
\|\Delta a_j\|\le\sum_{k\le n}\eta(1+\eta)^{k-1}\|a_j\|=\bigl((1+\eta)^n-1\bigr)\|a_j\|,
\]
and $(1+\eta)^n-1\le e^{n\eta}-1\le\frac{n\eta}{1-n\eta}$. The vector $b$ is
treated exactly like a column $j>n$, with $n$ steps of part (1).
\end{proof}

Note the elegance of the trick: we never compare the computed reflections with the exact ones of the problem, which may be ill-conditioned; we only use
that each computed vector $\hat v_k$ defines \emph{some} exact reflection, and that orthogonal maps do not enlarge perturbations.

\begin{corollary}[least squares by QR]\label{cor:lsq}
Let $A$ be $m\times n$ of rank $n$. Computing $\hat R$ and $\hat c=\widehat{Q^{\mathsf T}b}$ by Householder reflections and then solving $\hat R_1x=\hat c_1$ (the
top $n$ rows) by back substitution gives an $\hat x$ which minimises $\|(A+\delta A)x-(b+\delta b)\|$ exactly, with
$\|\delta A\|_F\le\varepsilon\|A\|_F$ and $\|\delta b\|\le\varepsilon\|b\|$, where $\varepsilon=\frac{n\eta}{1-n\eta}+\gamma_n\bigl(1+\frac{n\eta}{1-n\eta}\bigr)$.
\end{corollary}
\begin{proof}
By Theorem~\ref{thm:backsub}, $(\hat R_1+\Delta R)\hat x=\hat c_1$ with $\|\Delta R\|_F\le\gamma_n\|\hat R_1\|_F$, and
$\|\hat R\|_F=\|A+\Delta A\|_F\le(1+\frac{n\eta}{1-n\eta})\|A\|_F$. Let $\delta A=\Delta A+Q\begin{psmallmatrix}\Delta R\\0\end{psmallmatrix}$ and $\delta b=\Delta b$,
so that $A+\delta A=Q\begin{psmallmatrix}\hat R_1+\Delta R\\0\end{psmallmatrix}$ and $b+\delta b=Q\hat c$. For every $x$,
$\|(A+\delta A)x-(b+\delta b)\|^2=\|(\hat R_1+\Delta R)x-\hat c_1\|^2+\|\hat c_2\|^2$, which is minimal at $x=\hat x$.
\end{proof}

Compare with the \emph{normal equations} $A^{\mathsf T}Ax=A^{\mathsf T}b$ of Tome~S1, Theorem~4.8: they are mathematically equivalent, but
$\kappa(A^{\mathsf T}A)=\kappa(A)^2$ (Exercise~\ref{ex:condprops}), so forming $A^{\mathsf T}A$ squares the condition number and throws away half the digits
before any solving starts. When $\kappa(A)=10^8$, the normal equations in double precision can return nonsense, while $QR$ still gives about eight digits.

\begin{think}
The Gram--Schmidt process (Tome~S1, Theorem~4.5) also computes $A=QR$. Is it backward stable? Try to predict what happens to the orthogonality of the
computed columns $\hat q_j$ when the columns of $A$ are nearly dependent.
\end{think}

Classical Gram--Schmidt computes $q_j$ by subtracting from $a_j$ its projections $(q_i^{\mathsf T}a_j)q_i$ on the previous $q_i$. When $a_j$ is nearly in
the span of the $q_i$, the difference is small and dominated by rounding errors, which are not orthogonal to anything: the computed $\hat q_j$ can be far
from orthogonal to the earlier ones. For the matrix of Peter L\"auchli (1961) with columns $(1,\varepsilon,0,0)$, $(1,0,\varepsilon,0)$, $(1,0,0,\varepsilon)$ and
$\varepsilon^2<u$, classical Gram--Schmidt returns $\hat q_2^{\mathsf T}\hat q_3=\frac12$ (Exercise~\ref{ex:lauchli}). The \emph{modified} version, which projects
the current vector rather than the original $a_j$, keeps $\|\hat Q^{\mathsf T}\hat Q-I\|$ of order $u\kappa(A)$; Householder keeps it of order $u$, by
Theorem~\ref{thm:qrstable}. The notebook plots all three.

\section{Eigenvalues: from the power method to QR}\label{sec:eigen}

Eigenvalues are roots of the characteristic polynomial, and Abel and Galois showed there is no formula for roots of degree five or more: every
eigenvalue algorithm must be iterative. Worse, computing the characteristic polynomial and then its roots is a numerical disaster, since roots are
very sensitive to coefficients (Wilkinson's famous polynomial $\prod_{k=1}^{20}(x-k)$ loses its real roots when one coefficient changes by $2^{-23}$). The
algorithms work with the matrix itself.

\emph{The power method.} Let $A$ be diagonalisable with eigenvalues $|\lambda_1|>|\lambda_2|\ge\dots\ge|\lambda_n|$ and eigenvectors $x_1,\dots,x_n$. Start
from $q_0$ and iterate $q_{k}=Aq_{k-1}/\|Aq_{k-1}\|$.

\begin{proposition}[power method]\label{prop:power}
If $q_0=\sum_ic_ix_i$ with $c_1\ne0$, then the distance from $q_k$ to the line $\R x_1$ is $O(|\lambda_2/\lambda_1|^k)$. If moreover $A$ is symmetric, with
orthonormal $x_i$, then the Rayleigh quotient $\rho_k=q_k^{\mathsf T}Aq_k$ satisfies $|\rho_k-\lambda_1|\le2\|A\|\tan^2\phi_k$, where $\phi_k$ is the angle
between $q_k$ and $x_1$, with $\tan\phi_k\le\tan\phi_0\,|\lambda_2/\lambda_1|^k$.
\end{proposition}
\begin{proof}
$A^kq_0=\lambda_1^k\bigl(c_1x_1+\sum_{i\ge2}c_i(\lambda_i/\lambda_1)^kx_i\bigr)$, and $q_k$ is this vector normalised; the sum is $O(|\lambda_2/\lambda_1|^k)$.
In the symmetric case, write $q_k=\cos\phi_k\,x_1+\sin\phi_k\,w$ with $w\perp x_1$ a unit vector. The tangent of the angle is the ratio of the
norms of the components, $\tan\phi_k=\|\sum_{i\ge2}c_i\lambda_i^kx_i\|/|c_1\lambda_1^k|\le|\lambda_2/\lambda_1|^k\tan\phi_0$. Then
$\rho_k=\cos^2\phi_k\lambda_1+\sin^2\phi_k\,w^{\mathsf T}Aw$ (the cross terms vanish as $Ax_1=\lambda_1x_1\perp w$), so
$|\rho_k-\lambda_1|=\sin^2\phi_k|w^{\mathsf T}Aw-\lambda_1|\le2\|A\|\sin^2\phi_k\le2\|A\|\tan^2\phi_k$.
\end{proof}

For symmetric matrices the Rayleigh quotient doubles the number of correct digits of the vector: an eigenvector known to $8$ digits gives an eigenvalue
to $16$. Applied to $(A-\mu I)^{-1}$, the power method is \emph{inverse iteration}, which converges to the eigenvector of the eigenvalue nearest to $\mu$ at
rate $|\lambda-\mu|/|\lambda'-\mu|$, very fast when $\mu$ is a good guess (Exercise~\ref{ex:inverseiteration}). Taking for $\mu$ the current Rayleigh quotient
gives \emph{Rayleigh quotient iteration}, which for symmetric matrices converges cubically.

\emph{Subspaces and QR.} To find all eigenvalues, iterate a whole orthonormal basis: $Z_k=AQ_{k-1}$, then $Z_k=Q_kR_k$ (re-orthonormalise). This
\emph{simultaneous iteration} is the power method on each initial segment of columns at once. In 1961 John Francis in London, and independently Vera
Kublanovskaya in Leningrad, found a way to run it implicitly, on the matrix rather than on the basis. Heinz Rutishauser's LR algorithm of 1958 had shown the
way, with triangular factorisations.

\begin{definition}[the QR algorithm]\label{def:qralg}
Put $A_0=A$, and for $k\ge1$ factor $A_{k-1}=Q_kR_k$ (QR factorisation, with $R_k$ having positive diagonal) and set $A_k=R_kQ_k$.
\end{definition}

Then $A_k=Q_k^{\mathsf T}A_{k-1}Q_k$, so all the $A_k$ are orthogonally similar to $A$ and have its eigenvalues. Why should they converge to anything?

\begin{proposition}\label{prop:qrpowers}
Let $\mathcal Q_k=Q_1Q_2\cdots Q_k$ and $\mathcal R_k=R_kR_{k-1}\cdots R_1$. Then $A_k=\mathcal Q_k^{\mathsf T}A\mathcal Q_k$ and $A^k=\mathcal Q_k\mathcal R_k$:
$\mathcal Q_k\mathcal R_k$ is the $QR$ factorisation of the $k$th power of $A$.
\end{proposition}
\begin{proof}
The first claim follows from $A_k=Q_k^{\mathsf T}A_{k-1}Q_k$ by induction. For the second, by induction:
$\mathcal Q_k\mathcal R_k=\mathcal Q_{k-1}(Q_kR_k)\mathcal R_{k-1}=\mathcal Q_{k-1}A_{k-1}\mathcal R_{k-1}=A\mathcal Q_{k-1}\mathcal R_{k-1}=A\cdot A^{k-1}$,
using $\mathcal Q_{k-1}A_{k-1}=A\mathcal Q_{k-1}$.
\end{proof}

So the first $j$ columns of $\mathcal Q_k$ span $A^k$ applied to the first $j$ coordinate vectors: the $QR$ algorithm is simultaneous iteration started from
the identity, in disguise. The convergence theorem makes this precise in the symmetric case.

\begin{theorem}[convergence of the QR algorithm]\label{thm:qrconv}
Let $A=X\Lambda X^{\mathsf T}$ be symmetric with $X$ orthogonal and $\Lambda=\operatorname{diag}(\lambda_1,\dots,\lambda_n)$,
$|\lambda_1|>|\lambda_2|>\dots>|\lambda_n|>0$, and suppose all leading principal minors of $X^{\mathsf T}$ are nonzero. Then $A_k\to\Lambda$ up to the signs of
the off-diagonal entries; precisely, there are diagonal sign matrices $D_k$ with $D_kA_kD_k\to\Lambda$, and the entry $(A_k)_{ij}$, $i\ne j$, is
$O(\max_{l}|\lambda_{l+1}/\lambda_l|^k)$.
\end{theorem}
\begin{proof}
The hypothesis means $X^{\mathsf T}=LU$ with $L$ unit lower triangular and $U$ upper triangular and invertible (Gaussian elimination without row swaps
succeeds exactly when the leading minors are nonzero). Then
\[
A^k=X\Lambda^kLU=X(\Lambda^kL\Lambda^{-k})\Lambda^kU,
\]
and $\Lambda^kL\Lambda^{-k}=I+E_k$, where $(E_k)_{ij}=l_{ij}(\lambda_i/\lambda_j)^k$ for $i>j$ and $0$ otherwise, so $\|E_k\|=O(\rho^k)$ with
$\rho=\max_l|\lambda_{l+1}/\lambda_l|<1$. Let $I+E_k=\tilde Q_k\tilde R_k$ be the $QR$ factorisation with positive diagonal. The map sending an invertible
matrix to the factors of its $QR$ factorisation with positive diagonal is given by the Gram--Schmidt formulas, which are rational functions and square
roots of the entries, smooth near $I$, where it takes the value $(I,I)$; so $\tilde Q_k=I+O(\rho^k)$ and $\tilde R_k=I+O(\rho^k)$. Now
$A^k=(X\tilde Q_k)(\tilde R_k\Lambda^kU)$, an orthogonal matrix times an upper triangular one. By uniqueness of the $QR$ factorisation up to signs,
$\mathcal Q_k=X\tilde Q_kD_k$ for a diagonal sign matrix $D_k$ (chosen so that the diagonal of $D_k\tilde R_k\Lambda^kU$ is positive). By
Proposition~\ref{prop:qrpowers},
\[
D_kA_kD_k=D_k\mathcal Q_k^{\mathsf T}A\mathcal Q_kD_k=\tilde Q_k^{\mathsf T}X^{\mathsf T}AX\tilde Q_k=\tilde Q_k^{\mathsf T}\Lambda\tilde Q_k=\Lambda+O(\rho^k).
\]
\end{proof}

The hypothesis on the minors holds for almost every $A$, and rounding errors enforce it in practice. The off-diagonal entry in position $(l+1,l)$ decays
like $|\lambda_{l+1}/\lambda_l|^k$, which is slow when two eigenvalues are close (Figure~\ref{fig:qrconv}). Three ideas make the algorithm fast.

\begin{enumerate}
\item \emph{Hessenberg form.} A matrix is \emph{upper Hessenberg} if $a_{ij}=0$ for $i>j+1$. Householder reflections reduce any $A$ to Hessenberg form
$H=Q^{\mathsf T}AQ$ in $\frac{10}3n^3$ operations (Exercise~\ref{ex:hessenberg}); for symmetric $A$, $H$ is tridiagonal. A $QR$ step preserves the
Hessenberg form (Proposition~\ref{prop:hesspreserved}) and costs only $O(n^2)$, or $O(n)$ for tridiagonal matrices.
\item \emph{Shifts.} Factor $A_{k-1}-\mu_kI=Q_kR_k$ and set $A_k=R_kQ_k+\mu_kI$; this is again an orthogonal similarity, and the entry $(n,n-1)$ now decays like
$|(\lambda_n-\mu_k)/(\lambda_{n-1}-\mu_k)|$: with $\mu_k$ the current $(A_{k-1})_{nn}$, a Rayleigh quotient, convergence becomes cubic in the symmetric case. With
the \emph{Wilkinson shift}, the eigenvalue of the trailing $2\times2$ block closer to $(A_{k-1})_{nn}$, convergence is guaranteed for symmetric tridiagonal
matrices (Wilkinson, 1968). When $(A_k)_{n,n-1}$ is negligible, $\lambda_n$ is found and the problem \emph{deflates} to size $n-1$.
\item \emph{Implicit double shifts.} For real nonsymmetric matrices with complex eigenvalues, Francis showed how to apply two complex conjugate shifts at
once in real arithmetic, by chasing a small bulge down the Hessenberg matrix.
\end{enumerate}

With these, all eigenvalues of a symmetric $n\times n$ matrix cost about $\frac43n^3$ operations for the tridiagonal reduction and $O(n^2)$ for the
iteration, typically about two $QR$ steps per eigenvalue; and the whole computation is backward stable, being a sequence of orthogonal transformations.

\begin{proposition}\label{prop:hesspreserved}
If $H$ is upper Hessenberg and $H-\mu I=QR$ with $R$ invertible, then $RQ+\mu I$ is upper Hessenberg.
\end{proposition}
\begin{proof}
$Q=(H-\mu I)R^{-1}$, and $R^{-1}$ is upper triangular, so each column of $Q$ is a combination of the same and earlier columns of $H-\mu I$: $Q$ is upper
Hessenberg. Then $RQ$, an upper triangular matrix times an upper Hessenberg one, is upper Hessenberg: $(RQ)_{ij}=\sum_kr_{ik}q_{kj}$ needs $i\le k\le j+1$.
\end{proof}

\begin{figure}[ht]
\centering
\fig{ch01-qrconv}
\caption{The unshifted $QR$ algorithm on a symmetric $5\times5$ matrix with eigenvalues $10,7,6,3,1$: the subdiagonal entries $|(A_k)_{l+1,l}|$ decay
like $|\lambda_{l+1}/\lambda_l|^k$ (dashed lines), slowest for the close pair $7,6$. With Wilkinson shifts (dots), the last entry falls below $10^{-15}$
in five steps, the error being roughly cubed at each of the last ones.}
\label{fig:qrconv}
\end{figure}

Francis, then in his mid-twenties, left numerical analysis in 1962. When Gene Golub and Frank Uhlig tracked him down in 2007, he had no idea that his
algorithm had become one of the most used in scientific computing; he gave the opening talk at a meeting for its fiftieth anniversary in 2009.

\section{The singular value decomposition}\label{sec:svdcompute}

Tome~S1, Chapter~6 proves that every $m\times n$ matrix is $A=U\Sigma V^{\mathsf T}$, and explains what the singular values do: best low rank approximation,
the pseudo-inverse, principal components. How are they computed? The tempting route, eigenvalues of $A^{\mathsf T}A$, squares the condition number:
singular values below $\sqrt u\,\sigma_1$, that is $10^{-8}\sigma_1$, are lost in forming $A^{\mathsf T}A$ (Exercise~\ref{ex:squaring}). The good algorithms
never form it.

\emph{Golub and Kahan, 1965.} Alternating Householder reflections from the left (to annihilate a column below the diagonal) and from the right (to
annihilate a row to the right of the superdiagonal) reduces $A$ to upper \emph{bidiagonal} form.

\begin{proposition}[bidiagonalisation]\label{prop:bidiag}
For every $m\times n$ matrix $A$ with $m\ge n$ there are orthogonal $U,V$, products of at most $n$ and $n-2$ Householder reflections, with $U^{\mathsf T}AV=B$ upper
bidiagonal: $b_{ij}=0$ unless $j\in\{i,i+1\}$. The singular values of $A$ are those of $B$.
\end{proposition}
\begin{proof}
Step $k$ applies a reflection on the left acting on rows $k,\dots,m$ that zeroes column $k$ below the diagonal; it does not touch the earlier rows, so
the zeros made earlier stay. Then, if $k\le n-2$, a reflection on the right acting on columns $k+1,\dots,n$ zeroes row $k$ beyond position $k+1$; it does not
touch columns $1,\dots,k$, whose entries below row $k$ are already zero, so it does not destroy the zeros of column $k$. After $n$ steps the matrix is
upper bidiagonal. Orthogonal factors do not change singular values: $B^{\mathsf T}B=V^{\mathsf T}A^{\mathsf T}AV$ is similar to $A^{\mathsf T}A$.
\end{proof}

Then $B^{\mathsf T}B$ is symmetric tridiagonal, and the shifted $QR$ algorithm is applied to it \emph{implicitly}, working on $B$ alone by rotations,
so that $B^{\mathsf T}B$ is never formed (Golub and Kahan, 1965; Golub and Reinsch, 1970). This is what every library does, at a cost of
$O(mn^2)$ operations.

\emph{Jacobi, 1846.} A slower but strikingly simple algorithm goes back to Carl Gustav Jacobi, who used it in 1846, by hand, on the equations of the
secular perturbations of the planets. For a symmetric matrix $S$, let $\operatorname{off}(S)=\sum_{i\ne j}s_{ij}^2$. A \emph{Jacobi rotation} $J$ in the plane $(p,q)$,
by the angle $\theta$ with $\cot2\theta=(s_{qq}-s_{pp})/(2s_{pq})$, makes the entry $(p,q)$ of $J^{\mathsf T}SJ$ zero (Exercise~\ref{ex:jacobiangle}).

\begin{theorem}[convergence of Jacobi's method]\label{thm:jacobi}
Let $S$ be symmetric $n\times n$, $n\ge2$, and let $S'=J^{\mathsf T}SJ$ where $J$ is the Jacobi rotation annihilating an entry $(p,q)$ of maximal modulus
among the off-diagonal entries. Then $\operatorname{off}(S')\le\bigl(1-\frac2{n(n-1)}\bigr)\operatorname{off}(S)$. In particular the iterates converge to a
diagonal matrix carrying the eigenvalues of $S$.
\end{theorem}
\begin{proof}
The rotation changes only rows and columns $p$ and $q$. The Frobenius norm is invariant, $\|S'\|_F=\|S\|_F$, and so is the Frobenius norm of the
$2\times2$ block on rows and columns $p,q$, which is rotated: $s'^2_{pp}+s'^2_{qq}+2s'^2_{pq}=s_{pp}^2+s_{qq}^2+2s_{pq}^2$. For $i\notin\{p,q\}$, the pair
$(s_{ip},s_{iq})$ is rotated, so $s'^2_{ip}+s'^2_{iq}=s_{ip}^2+s_{iq}^2$, and other entries are unchanged. Hence the sum of the squares of the
off-diagonal entries outside the block is unchanged, while $s'_{pq}=0$: $\operatorname{off}(S')=\operatorname{off}(S)-2s_{pq}^2$. As $s_{pq}^2$ is the
largest of the $n(n-1)$ terms of $\operatorname{off}(S)$, $2s_{pq}^2\ge\frac2{n(n-1)}\operatorname{off}(S)$. So $\operatorname{off}\to0$ geometrically, and
the diagonal parts $D_k$ satisfy $\|S_k-D_k\|_F\to0$ where $S_k$ is similar to $S$. The eigenvalues of $S_k$ and those of $D_k$ (its diagonal entries, sorted)
then differ by at most $\|S_k-D_k\|$ (Weyl's inequality, Tome~S1, Corollary~5.7), so the sorted diagonal converges to the sorted eigenvalues.
\end{proof}

The theorem gives only linear convergence, but in fact Jacobi's method becomes quadratically convergent once off$(S)$ is small (Sch\"onhage, 1961), and it
is often \emph{more} accurate than $QR$ for the small eigenvalues of well-scaled matrices (Demmel and Veseli\'c, 1992). For the SVD, the \emph{one-sided}
version of Magnus Hestenes (1958) applies the rotations to the columns of $A$: rotating columns $p,q$ of $A$ by the Jacobi rotation of $A^{\mathsf T}A$
for $(p,q)$ makes them orthogonal, and the three numbers needed, $\|a_p\|^2$, $\|a_q\|^2$ and $a_p^{\mathsf T}a_q$, are computed from the columns at each step
rather than from a stored $A^{\mathsf T}A$. When all columns are orthogonal, $AV=U\Sigma$ with $\sigma_j$ the column norms. The notebook builds an SVD this
way in twenty lines and compares it with the library.

\begin{beauty}
The $QR$ algorithm hides a power method inside a similarity: by Proposition~\ref{prop:qrpowers}, factoring and multiplying back in reverse order
computes, without ever forming it, the orthonormal basis of $A^k$ applied to the coordinate axes. Each eigenvalue emerges at the corner where it has
won its race against the others, at the speed $|\lambda_{l+1}/\lambda_l|$ dictated by the power method, and the shifts turn the race into a sprint.
\end{beauty}

\section{Exercises}

\begin{exercise}[$\star$]\label{ex:summation}
Show that recursive summation of $x_1,\dots,x_n\ge0$ gives $|\mathrm{fl}(\sum x_i)-\sum x_i|\le\gamma_{n-1}\sum x_i$, and that the error bound is smaller
when the terms are added in increasing order. Compute $\sum_{k=1}^{N}1/k^2$ in single precision both ways for $N=10^6$, and explain.
\end{exercise}

\begin{exercise}[$\star$]\label{ex:quadratic}
The roots of $x^2-2bx+c$ with $b^2>c>0$ and $b>0$ are $b\pm\sqrt{b^2-c}$. Explain why the formula for the smaller root loses accuracy when $c\ll b^2$, and
give a formula for it without cancellation.
\end{exercise}

\begin{exercise}[$\star$]\label{ex:condprops}
Show that $\kappa(AB)\le\kappa(A)\kappa(B)$, that $\kappa(Q)=1$ for orthogonal $Q$, and that $\kappa(A^{\mathsf T}A)=\kappa(A)^2$ for $A$ of full column rank
(with $\kappa=\sigma_1/\sigma_n$ for rectangular matrices).
\end{exercise}

\begin{exercise}[$\star$]\label{ex:nopivot}
For $A=\begin{psmallmatrix}\varepsilon&1\\1&1\end{psmallmatrix}$ with $\varepsilon=10^{-20}$ and $b=(1,2)$, follow Gaussian elimination without pivoting in
double precision and find the computed solution. Then do it with partial pivoting. Which part of Theorem~\ref{thm:geerror} explains the failure?
\end{exercise}

\begin{exercise}[$\star$]\label{ex:wilkinsongrowth}
Prove that partial pivoting performs no row swaps on the matrix $W_n$ of Proposition~\ref{prop:growth}, and that $u_{nn}=2^{n-1}$. Show also that
$W_n$ is well conditioned: $\kappa_\infty(W_n)=\|W_n\|_\infty\|W_n^{-1}\|_\infty=n$.
\end{exercise}

\begin{exercise}[$\star\star$]\label{ex:lauchli}
Let $\varepsilon>0$ with $\mathrm{fl}(1+\varepsilon^2)=1$. Run classical Gram--Schmidt on the columns $a_1=(1,\varepsilon,0,0)$, $a_2=(1,0,\varepsilon,0)$,
$a_3=(1,0,0,\varepsilon)$, assuming that the only rounding errors are $\mathrm{fl}(1+\varepsilon^2)=1$, and show that $\hat q_2^{\mathsf T}\hat q_3=\frac12$.
Then do the same with modified Gram--Schmidt, where $a_3$ is first projected off $\hat q_1$ and the \emph{result} is projected off $\hat q_2$.
\end{exercise}

\begin{exercise}[$\star\star$]\label{ex:hilbertcond}
(a) Show that the Hilbert matrix is the Gram matrix of the functions $x^{i-1}$ in $L^2(0,1)$, hence positive definite. (b) Show that $\|H_n^{-1}\|\ge
1/\min_p\int_0^1p^2$, the minimum over polynomials of degree $<n$ with coefficient vector of norm $1$, and use the shifted Legendre polynomial to show that
$\kappa(H_n)$ grows at least like $c^n$ for some $c>1$.
\end{exercise}

\begin{exercise}[$\star$]\label{ex:inverseiteration}
Let $A$ be diagonalisable with eigenvalues $\lambda_i$, and $\mu$ not an eigenvalue, with $\lambda$ the unique eigenvalue nearest to $\mu$ and $\lambda'$ the
next nearest. Show that the power method applied to $(A-\mu I)^{-1}$ converges to an eigenvector for $\lambda$ at rate $|\lambda-\mu|/|\lambda'-\mu|$. Why is the
near singularity of $A-\mu I$ harmless?
\end{exercise}

\begin{exercise}[$\star\star$]\label{ex:hessenberg}
Show how $n-2$ Householder reflections, applied on both sides, reduce an $n\times n$ matrix to upper Hessenberg form by an orthogonal similarity, and
that the cost is about $\frac{10}3n^3$ operations. Why can one not reduce all the way to triangular form in the same way?
\end{exercise}

\begin{exercise}[$\star$]\label{ex:jacobiangle}
Let $S$ be symmetric and $J$ the rotation by $\theta$ in the plane of coordinates $p<q$. Compute the entry $(p,q)$ of $J^{\mathsf T}SJ$ and show it vanishes
when $\cot2\theta=(s_{qq}-s_{pp})/(2s_{pq})$. Show that one can always choose $|\theta|\le\pi/4$.
\end{exercise}

\begin{exercise}[$\star\star$]\label{ex:squaring}
Let $A=\begin{psmallmatrix}1&1\\\delta&0\\0&\delta\end{psmallmatrix}$ with $\delta^2<u$. Compute its singular values exactly, then show that the
floating-point $A^{\mathsf T}A$ is singular. What does the one-sided Jacobi method compute instead?
\end{exercise}

\begin{exercise}[$\star\star\star$]\label{ex:givens}
A \emph{Givens rotation} acts on two coordinates. Show that $QR$ factorisation by Givens rotations, zeroing the entries below the diagonal one at a
time, satisfies the conclusion of Theorem~\ref{thm:qrstable} with a constant of order $m+n$, by proving the analogue of Lemma~\ref{lem:onereflection} for
one rotation and counting how many rotations touch each column. When are rotations preferable to reflections?
\end{exercise}

\section{Hints}

\begin{hint}{ex:summation}
Use Lemma~\ref{lem:oneeq} with all $b$'s equal to $1$; the term added first is rounded the most times.
\end{hint}
\begin{hint}{ex:quadratic}
The product of the roots is $c$.
\end{hint}
\begin{hint}{ex:condprops}
$\|AB\|\le\|A\|\|B\|$; for the last one, use the SVD of $A$.
\end{hint}
\begin{hint}{ex:nopivot}
Compute $u_{22}=1-10^{20}$ in floating point, then back substitute.
\end{hint}
\begin{hint}{ex:wilkinsongrowth}
At step $k$, all entries of column $k$ on and below the diagonal are $\pm1$. For the inverse, solve $W_nx=b$ using the partial sums $x_1+\dots+x_i$, or guess the inverse for $n=3,4$ and check it.
\end{hint}
\begin{hint}{ex:lauchli}
$\hat q_1=(1,\varepsilon,0,0)$ after rounding; then $\hat q_2=(0,-1,1,0)/\sqrt2$.
\end{hint}
\begin{hint}{ex:hilbertcond}
$x^{\mathsf T}H_nx=\int_0^1(\sum_ix_it^{i-1})^2\dd t$. The shifted Legendre polynomial of degree $n-1$ has integer coefficients of size about $(3+2\sqrt2)^n$.
\end{hint}
\begin{hint}{ex:inverseiteration}
The eigenvalues of $(A-\mu I)^{-1}$ are $1/(\lambda_i-\mu)$. The error in solving with $A-\mu I$ is mostly along the wanted eigenvector.
\end{hint}
\begin{hint}{ex:hessenberg}
At step $k$ the reflection acts on rows $k+1,\dots,n$, so applying it on the right does not touch column $k$.
\end{hint}
\begin{hint}{ex:jacobiangle}
The $(p,q)$ entry is $\cos2\theta\,s_{pq}+\tfrac12\sin2\theta\,(s_{pp}-s_{qq})$.
\end{hint}
\begin{hint}{ex:squaring}
$A^{\mathsf T}A=\begin{psmallmatrix}1+\delta^2&1\\1&1+\delta^2\end{psmallmatrix}$.
\end{hint}
\begin{hint}{ex:givens}
A rotation touches two entries; each entry of a column is touched at most $m+n$ times in a suitable order.
\end{hint}

\section{Solutions}

\begin{solution}{ex:summation}
\idea Summation is the special case of Lemma~\ref{lem:oneeq} with no products and no division; the error depends on which terms are rounded most often.
\writeup Recursive summation $s_1=x_1$, $s_i=\mathrm{fl}(s_{i-1}+x_i)$ gives $\hat s_n=\sum_ix_i(1+\theta_i)$ where $x_1$ and $x_2$ are rounded $n-1$ times and $x_i$
($i\ge3$) $n-i+1$ times, so $|\theta_i|\le\gamma_{n-1}$ and, all terms being nonnegative, $|\hat s_n-\sum x_i|\le\gamma_{n-1}\sum x_i$. A sharper bound is
$\sum_i|\theta_i|x_i\lesssim u\sum_i(n-i+1)x_i$ (for $i\ge2$), which is smallest when the large weights $n-i+1$ multiply the small terms: add in increasing
order. Equivalently, the error made at step $i$ is at most $u\hat s_i$, and the sum of the partial sums is smallest when the small terms come first. For
$\sum_{k\le10^6}1/k^2$ in single precision ($u=2^{-24}$), adding from $k=1$ the partial sum reaches $1.64$ quickly, and later terms below $u\cdot1.64\approx10^{-7}$,
that is for $k$ beyond a few thousand, are entirely lost: the sum stalls at $1.6447\dots$ with an error of about $2\times10^{-4}$. Adding from $k=10^6$ down, the
small terms accumulate first and the result is correct to about seven digits, $1.644933$ (the exact value is $\pi^2/6-\sum_{k>10^6}k^{-2}\approx1.6449331$).
The notebook checks these numbers.
\end{solution}

\begin{solution}{ex:quadratic}
\idea Cancellation is harmless in itself; it is the subtraction of two nearly equal computed numbers that loses digits. Avoid it by an algebraic identity.
\writeup When $c\ll b^2$, $\sqrt{b^2-c}\approx b$, and $b-\sqrt{b^2-c}$ subtracts two nearly equal numbers, the second of which carries a rounding error of size
$ub$: the relative error of the difference, which is about $c/(2b)$, is of order $ub/(c/2b)=2ub^2/c$, large when $c\ll b^2$. Compute instead the large root
$x_+=b+\sqrt{b^2-c}$, a sum of positive numbers with relative error a few $u$, and then $x_-=c/x_+$, since $x_+x_-=c$. This also has relative error a few $u$.
\end{solution}

\begin{solution}{ex:condprops}
\idea Write everything in terms of norms and singular values.
\writeup
\[
\kappa(AB)=\|AB\|\|B^{-1}A^{-1}\|\le\|A\|\|B\|\|B^{-1}\|\|A^{-1}\|=\kappa(A)\kappa(B).
\]
For orthogonal $Q$, $\|Q\|=\|Q^{-1}\|=1$. If $A=U\Sigma V^{\mathsf T}$
(thin SVD, $\Sigma$ of size $n$), then $A^{\mathsf T}A=V\Sigma^2V^{\mathsf T}$ has eigenvalues, hence singular values, $\sigma_i^2$, so
$\kappa(A^{\mathsf T}A)=\sigma_1^2/\sigma_n^2=\kappa(A)^2$.
\end{solution}

\begin{solution}{ex:nopivot}
\idea Find the single operation where information is lost.
\writeup Without pivoting $l_{21}=1/\varepsilon=10^{20}$, $u_{22}=\mathrm{fl}(1-10^{20})=-10^{20}$ (the $1$ is far below the last digit), and
$y_2=\mathrm{fl}(2-10^{20}\cdot1)=-10^{20}$. Back substitution: $x_2=y_2/u_{22}=1$, then $x_1=(1-x_2)/\varepsilon=0$. The true solution is
$x_1=1/(1-\varepsilon)\approx1$, $x_2=(1-2\varepsilon)/(1-\varepsilon)\approx1$; the computed one has $x_1$ completely wrong although $\kappa(A)\approx2.6$. In
Theorem~\ref{thm:geerror}, $|\hat L||\hat U|=\begin{psmallmatrix}\varepsilon&1\\1&2\cdot10^{20}\end{psmallmatrix}$ is huge compared with $|A|$: the backward
error bound is useless. With partial pivoting the rows are swapped, $l_{21}=\varepsilon$, $u_{22}=\mathrm{fl}(1-\varepsilon)=1$, and the computed solution is
$x_2=\mathrm{fl}(1-2\varepsilon)=1$, $x_1=2-x_2=1$, correct to full precision.
\end{solution}

\begin{solution}{ex:wilkinsongrowth}
\idea Follow the first two steps by hand and guess the invariant.
\writeup Claim: before step $k$, the active submatrix (rows and columns $\ge k$) has $1$ on its diagonal, $-1$ below it, $0$ above it except in the last column,
whose entries are $2^{k-1}$ in rows $k,\dots,n$. True for $k=1$. In column $k$ all candidates for the pivot have modulus $1$, so (with ties broken in favour
of the current row) no swap occurs, $l_{ik}=-1$ for $i>k$, and step $k$ adds row $k$ to each later row: the entries in columns $k+1,\dots,n-1$ of row $k$ are $0$,
so these columns are unchanged, and the last column becomes $2^{k-1}+2^{k-1}=2^k$ in rows $k+1,\dots,n$. At the end $u_{nn}=2^{n-1}$ while $\max|a_{ij}|=1$, so
$\rho_n=2^{n-1}$. For the conditioning, $\|W_n\|_\infty=n$ (the last row). The inverse is the matrix $V$ with, for $i<n$, $v_{ii}=\frac12$,
$v_{ij}=-2^{-(j-i+1)}$ for $i<j<n$, $v_{in}=-2^{-(n-i)}$, zeros below the diagonal, and last row $v_{nj}=2^{-j}$ for $j<n$, $v_{nn}=2^{1-n}$; one checks $VW_n=I$ by
multiplying out (the notebook does it for $n\le8$; this formula is found by solving $W_nx=b$ with the partial sums $S_i=x_1+\dots+x_i$, which satisfy
$S_i=2S_{i-1}+b_i-x_n$). Each row of $V$ has absolute sum $\frac12+(\frac12-2^{-(n-i)})+2^{-(n-i)}=1$ for $i<n$, and $\sum_{j<n}2^{-j}+2^{1-n}=1$ for the last.
So $\|W_n^{-1}\|_\infty=1$ and $\kappa_\infty(W_n)=n$: the growth $2^{n-1}$ has nothing to do with ill conditioning, and the computed solution of $W_nx=b$ can
be badly wrong for $n=60$ although the problem is perfectly well posed.
\end{solution}

\begin{solution}{ex:lauchli}
\idea Only $1+\varepsilon^2$ is rounded, so everything else can be followed exactly; watch which vector is projected.
\writeup $\|a_1\|=\sqrt{1+\varepsilon^2}$ rounds to $1$, so $\hat q_1=a_1$. Then $\hat q_1^{\mathsf T}a_2=1$ (rounded), $a_2-\hat q_1=(0,-\varepsilon,\varepsilon,0)$, and
$\hat q_2=(0,-1,1,0)/\sqrt2$. Classical: $\hat q_1^{\mathsf T}a_3=1$ and $\hat q_2^{\mathsf T}a_3=0$, so $a_3-\hat q_1=(0,-\varepsilon,0,\varepsilon)$ and
$\hat q_3=(0,-1,0,1)/\sqrt2$. Then $\hat q_2^{\mathsf T}\hat q_3=\frac12$: the computed vectors make an angle of $60^\circ$. Modified: first
$w=a_3-(\hat q_1^{\mathsf T}a_3)\hat q_1=(0,-\varepsilon,0,\varepsilon)$, then $w-(\hat q_2^{\mathsf T}w)\hat q_2$ with $\hat q_2^{\mathsf T}w=\varepsilon/\sqrt2$, giving
$(0,-\frac\varepsilon2,-\frac\varepsilon2,\varepsilon)$ and $\hat q_3=(0,-1,-1,2)/\sqrt6$, now orthogonal to $\hat q_2$; and
$\hat q_1^{\mathsf T}\hat q_3=-\varepsilon/\sqrt6$, $\hat q_1^{\mathsf T}\hat q_2=-\varepsilon/\sqrt2$, of order $\varepsilon\approx\sqrt u$, consistent with $u\kappa(A)$
since $\kappa(A)\approx\sqrt3/\varepsilon$.
\end{solution}

\begin{solution}{ex:hilbertcond}
\idea A Gram matrix of nearly dependent functions is nearly singular; the monomials on $[0,1]$ are nearly dependent, and orthogonal polynomials measure how much.
\writeup (a) $\int_0^1t^{i-1}t^{j-1}\dd t=1/(i+j-1)$, so $x^{\mathsf T}H_nx=\int_0^1p_x(t)^2\dd t$ with $p_x(t)=\sum_ix_it^{i-1}$, positive unless $x=0$.
(b) $1/\|H_n^{-1}\|=\lambda_{\min}(H_n)=\min_{\|x\|=1}\int_0^1p_x^2$. Take $p$ the shifted Legendre polynomial $\tilde P_{n-1}(t)=P_{n-1}(2t-1)$, with
$\int_0^1\tilde P_{n-1}^2=1/(2n-1)$ and coefficient vector $c$; its leading coefficient is $\binom{2n-2}{n-1}\ge4^{n-1}/(2n-1)$, so $\|c\|\ge4^{n-1}/(2n-1)$ and
$\lambda_{\min}\le\frac{1/(2n-1)}{\|c\|^2}\le(2n-1)16^{-(n-1)}$. As $\|H_n\|\ge h_{11}=1$, $\kappa(H_n)\ge16^{n-1}/(2n-1)$. The true growth rate is about
$(1+\sqrt2)^{4n}\approx e^{3.53n}$, from the full coefficient vector rather than its leading term.
\end{solution}

\begin{solution}{ex:inverseiteration}
\idea Inverse iteration is the power method for a matrix whose dominant eigenvalue is the one we want.
\writeup $(A-\mu I)^{-1}$ has the same eigenvectors as $A$ and eigenvalues $1/(\lambda_i-\mu)$, the largest in modulus being $1/(\lambda-\mu)$ and the next
$1/(\lambda'-\mu)$. By Proposition~\ref{prop:power} the iterates converge to the line of the eigenvector for $\lambda$ at rate
$\frac{1/|\lambda'-\mu|}{1/|\lambda-\mu|}=\frac{|\lambda-\mu|}{|\lambda'-\mu|}$. Solving $(A-\mu I)z=q$ with a backward stable method gives the exact $z$ for
$A-\mu I+\Delta$, and the error, though large, is amplified by $(A-\mu I)^{-1}$ mainly along its dominant eigenvector, which is the direction we want; after
normalisation the computed vector is accurate. (Wilkinson made this precise; the notebook shows a shift within $10^{-14}$ of the eigenvalue converging in one step.)
\end{solution}

\begin{solution}{ex:hessenberg}
\idea Copy Householder $QR$, but leave one subdiagonal so that the right multiplication does not undo the work of the left one.
\writeup At step $k=1,\dots,n-2$, let $P_k$ act on coordinates $k+1,\dots,n$ and zero the entries $(k+2,k),\dots,(n,k)$ of column $k$. Then $P_kAP_k$: the
left multiplication changes rows $k+1,\dots,n$, creating the zeros; the right multiplication changes columns $k+1,\dots,n$ and so does not touch column
$k$, keeping them. Earlier columns are untouched by both since their entries in rows $k+1,\dots$ beyond the subdiagonal are already zero and the left
action mixes zeros. The left update on the trailing $(n-k)\times(n-k+1)$ block costs about $4(n-k)^2$ operations, and the right update on the $n\times(n-k)$
block about $4n(n-k)$; summing, $\sum_k(4(n-k)^2+4n(n-k))\approx\frac43n^3+2n^3=\frac{10}3n^3$. To reach triangular form we would have to zero column $1$
below the diagonal by a reflection acting on rows $1,\dots,n$; the right multiplication by it would then act on columns $1,\dots,n$ and fill column $1$ again.
Indeed no finite procedure can triangularise all matrices, since that would compute eigenvalues exactly in finitely many steps.
\end{solution}

\begin{solution}{ex:jacobiangle}
\idea Only a $2\times2$ problem is involved.
\writeup With $c=\cos\theta$, $s=\sin\theta$, the entry is $(J^{\mathsf T}SJ)_{pq}=cs(s_{pp}-s_{qq})+(c^2-s^2)s_{pq}=\frac12\sin2\theta(s_{pp}-s_{qq})+\cos2\theta\,s_{pq}$
for $J$ with $J_{pp}=J_{qq}=c$, $J_{pq}=s$, $J_{qp}=-s$. It vanishes iff $\cot2\theta=(s_{qq}-s_{pp})/(2s_{pq})$ (for $s_{pq}\ne0$). As $\cot$ takes every real value
on $(0,\pi)$, there is $2\theta\in(0,\pi)$, and replacing $2\theta$ by $2\theta-\pi$ if needed, $2\theta\in(-\pi/2,\pi/2]$, so $|\theta|\le\pi/4$. The small angle
keeps the rotation close to the identity, which is important for the asymptotic quadratic convergence.
\end{solution}

\begin{solution}{ex:squaring}
\idea Rounding kills the information in $A^{\mathsf T}A$, not in $A$.
\writeup $A^{\mathsf T}A=\begin{psmallmatrix}1+\delta^2&1\\1&1+\delta^2\end{psmallmatrix}$ has eigenvalues $2+\delta^2$ and $\delta^2$, so $\sigma_1=\sqrt{2+\delta^2}$ and
$\sigma_2=\delta$. In floating point $\mathrm{fl}(1+\delta^2)=1$, so the computed $A^{\mathsf T}A$ is $\begin{psmallmatrix}1&1\\1&1\end{psmallmatrix}$, singular, and the
computed $\sigma_2$ is $0$ (or a rounding error), with no correct digit. One-sided Jacobi rotates the columns $a_1=(1,\delta,0)$, $a_2=(1,0,\delta)$ by $45^\circ$
(as $\|a_1\|=\|a_2\|$), giving $(a_1+a_2)/\sqrt2=(\sqrt2,\delta/\sqrt2,\delta/\sqrt2)$ and $(a_1-a_2)/\sqrt2=(0,\delta/\sqrt2,-\delta/\sqrt2)$, which are computed with small
relative error in each entry; the second has norm $\delta$, correct to full relative precision.
\end{solution}

\begin{solution}{ex:givens}
\idea The proof of Theorem~\ref{thm:qrstable} only used that each step is an exact orthogonal map applied to a slightly perturbed column; this holds for rotations
too, and only the number of steps per column changes.
\writeup A rotation in the plane $(i,j)$ is defined by $c=a/r$, $s=b/r$ with $r=\sqrt{a^2+b^2}$; the computed $\hat c,\hat s$ have relative errors at most
$\gamma_4$, and define an exact rotation $\tilde G$ after renormalisation by $\sqrt{\hat c^2+\hat s^2}=1+O(\gamma_4)$. Applying $\hat G$ to the pair $(x_i,x_j)$ of a
vector gives $(\hat cx_i+\hat sx_j)(1+\theta)$ and similarly, so $\mathrm{fl}(\hat Gx)=\tilde G(x+\Delta x)$ with $\|\Delta x\|\le\gamma_{c}\|(x_i,x_j)\|\le\gamma_c\|x\|$
for a small absolute constant $c$. Zeroing the subdiagonal column by column, from the bottom up, each column is touched by at most $m-1$ rotations for each of the
$n$ columns processed, but the rotations in different rows are disjoint and can be grouped into at most $m+n-2$ stages, each stage an orthogonal map acting on
disjoint pairs, with error $\gamma_c\|x\|$ per stage. The telescoping argument of Theorem~\ref{thm:qrstable} then gives $\|\Delta a_j\|\le((1+\gamma_c)^{m+n-2}-1)\|a_j\|$.
Rotations cost about twice as much as reflections for dense matrices, but they touch two rows at a time, which is better for matrices that are already almost
triangular (a Hessenberg matrix needs only $n-1$ rotations), and for updating a factorisation when a row is added.
\end{solution}

\section*{Sources and further reading}

\begin{itemize}
\item \emph{A research paper to read.} A.~M.~Turing, \emph{Rounding-off errors in matrix processes}, Quarterly Journal of Mechanics and Applied Mathematics 1
(1948), 287--308. Read Section~1 for the setting, then the definition of the ``$N$-condition number'' and ``$M$-condition number'' and how they control the error of
the solution (our Theorem~\ref{thm:perturb}); Turing also introduces the triangular resolution $A=LDU$. Then read the opening sections of J.~H.~Wilkinson,
\emph{Error analysis of direct methods of matrix inversion}, Journal of the ACM 8 (1961), 281--330, where backward error analysis appears in full, with the growth
factor of Proposition~\ref{prop:growth}.
\item \emph{The stories.} J.~H.~Wilkinson, \emph{Some comments from a numerical analyst} (1970 Turing Award lecture), Journal of the ACM 18 (1971), 137--147, for
the twelve and the eighteen equations, the pessimism of 1946, Turing's defeat and his ``second round'', and all the quotations of our opening. The bounds growing like
$4^n$ are those of H.~Hotelling, \emph{Some new methods in matrix calculation}, Annals of Mathematical Statistics 14 (1943), 1--34. J.~von~Neumann and H.~H.~Goldstine,
\emph{Numerical inverting of matrices of high order}, Bulletin of the AMS 53 (1947), 1021--1099. The Patriot failure: United States General Accounting Office,
\emph{Patriot Missile Defense: Software Problem Led to System Failure at Dhahran, Saudi Arabia}, report GAO/IMTEC-92-26, February 1992 (the date, the 28 deaths, the
24-bit register, the 100 hours, the $0.3433$ seconds and the 687 metres of its Appendix~II, and the corrected software arriving on 26 February). G.~H.~Golub and F.~Uhlig,
\emph{The QR algorithm: 50 years later, its genesis by John Francis and Vera Kublanovskaya and subsequent developments}, IMA Journal of Numerical Analysis 29 (2009), 467--485,
for Francis' story.
\item A.~S.~Householder, \emph{Unitary triangularization of a nonsymmetric matrix}, Journal of the ACM 5 (1958), 339--342; J.~G.~F.~Francis, \emph{The QR transformation:
a unitary analogue to the LR transformation}, Computer Journal 4 (1961), 265--271 and 332--345; V.~N.~Kublanovskaya, \emph{On some algorithms for the solution of the
complete eigenvalue problem}, USSR Computational Mathematics and Mathematical Physics 1 (1962), 637--657; H.~Rutishauser, \emph{Solution of eigenvalue problems with the
LR-transformation}, National Bureau of Standards Applied Mathematics Series 49 (1958), 47--81; G.~H.~Golub and W.~Kahan, \emph{Calculating the singular values and
pseudo-inverse of a matrix}, SIAM Journal on Numerical Analysis B 2 (1965), 205--224; G.~H.~Golub and C.~Reinsch, \emph{Singular value decomposition and least squares
solutions}, Numerische Mathematik 14 (1970), 403--420; C.~G.~J.~Jacobi, \emph{\"Uber ein leichtes Verfahren die in der Theorie der S\"acularst\"orungen vorkommenden
Gleichungen numerisch aufzul\"osen}, Journal f\"ur die reine und angewandte Mathematik 30 (1846), 51--94; M.~R.~Hestenes, \emph{Inversion of matrices by
biorthogonalization and related results}, Journal of the SIAM 6 (1958), 51--90; P.~L\"auchli, \emph{Jordan-Elimination und Ausgleichung nach kleinsten Quadraten},
Numerische Mathematik 3 (1961), 226--240; A.~Sch\"onhage, \emph{Zur Konvergenz des Jacobi-Verfahrens}, Numerische Mathematik 3 (1961), 374--380; J.~H.~Wilkinson, \emph{Global convergence of tridiagonal QR algorithm with origin shifts}, Linear Algebra and its Applications 1 (1968), 409--420; J.~Demmel and K.~Veseli\'c, \emph{Jacobi's method is more accurate than QR}, SIAM Journal on Matrix Analysis and Applications 13
(1992), 1204--1245.
\item L.~N.~Trefethen and R.~S.~Schreiber, \emph{Average-case stability of Gaussian elimination}, SIAM Journal on Matrix Analysis and Applications 11 (1990), 335--360;
A.~Sankar, D.~A.~Spielman and S.-H.~Teng, \emph{Smoothed analysis of the condition numbers and growth factors of matrices}, SIAM Journal on Matrix Analysis and
Applications 28 (2006), 446--476.
\item N.~J.~Higham, \emph{Accuracy and Stability of Numerical Algorithms}, 2nd edition, SIAM, 2002, is the reference for Sections~\ref{sec:float}--\ref{sec:householder} (its
Chapters~2, 3, 8, 9 and 19; our constants are cruder than its). L.~N.~Trefethen and D.~Bau, \emph{Numerical Linear Algebra}, SIAM, 1997, is the most readable
introduction, and G.~H.~Golub and C.~F.~Van~Loan, \emph{Matrix Computations}, 4th edition, Johns Hopkins, 2013, the encyclopedia.
\end{itemize}
